Пусть $F \subset K$ -- расширение полей. 
Множество автоморфизмов $K$, оставляющих $F$ на месте является группой и
называется группой автоморфизмов и обозначается 
$Aut_F (K)$ = $Aut([K : F])$. 
Если $F$ -- основное поле ($\Q$ или $\Z_p$), то символ $F$ опускают.

а) $Aut_F (K)$ -- группа.

Покажем, что $Aut_F (K)$ является подгруппой $Aut\ K$.
Для этого нужно показать:

1) замкнутость

$$
\forall \varphi, \psi \in Aut_F (K)\ \varphi \circ \psi \in Aut_F (K)
$$

Это действительно так, так как 
$$
\forall x \in F \ (\varphi \circ \psi)(x) = x
$$

2) замкнутость относительно обратного

$$
\forall \varphi \in Aut_F (K)\ \forall x \in F\ \ 
\varphi^{-1}(x) = x
$$

б) Пусть $H \subset Aut_F (K)$ -- подгруппа. 
Тогда $K^H = \{x \in K\ |\ \forall h \in H \ h(x) = x\}$ 
является полем, причём $K \supset K^H \supset F$.

Проверим все аксиомы поля:

1) Замкнутость. Пусть $h(\alpha) = \alpha, h(\beta) = \beta$

$$
h(\alpha + \beta) = h(\alpha) + h(\beta) = \alpha + \beta
\Rightarrow \alpha + \beta \in K^H
$$

$$
h(\alpha \cdot \beta) = h(\alpha) \cdot h(\beta) = \alpha \cdot \beta
\Rightarrow \alpha \cdot \beta \in K^H
$$

2) Существование нейтрального. 

У любого автоморфизма $h(0) = 0, h(1) = 1$, поэтому $0, 1 \in K^H$.

3) Существование обратного

$$
h(-\alpha) = -\alpha, h(\alpha^{-1}) = \alpha^{-1} (\alpha \neq 0)
$$

4) Коммутативность. Следует из коммутативности в $K$.

5) Дистрибутивность. Следует из дистрибутивности в $K$.

$K \supset K^H$ очевидно.

$K^H \supset F$, так как $H \subset Aut_F (K)$, то есть 
точки из $F$ точно остаются неподвижными.
